\chapter{2D Truss Element}
\label{ch:truss}

\section{2D Truss Element}

A truss element can carry axial loads only. In a 2D structure, each node has
two translational DOFs: $u$ (horizontal) and $v$ (vertical). The element is
inclined at angle $\theta$ with respect to the global $x$-axis~\cite{cook2002concepts}.

\subsection{Local Stiffness Matrix}

In the local coordinate system aligned with the element axis, the stiffness
matrix is identical to the 1D bar element~\cref{eq:bar_stiffness}, expanded
to account for the transverse DOF (which has zero stiffness):
\begin{equation}
    \bar{\matr{k}}^e = \frac{EA}{L}
    \begin{bmatrix} 1 & 0 & -1 & 0 \\ 0 & 0 & 0 & 0 \\
                   -1 & 0 &  1 & 0 \\ 0 & 0 & 0 & 0 \end{bmatrix}
    \label{eq:truss_local}
\end{equation}
where the DOF ordering is $[\bar{u}_1,\, \bar{v}_1,\, \bar{u}_2,\, \bar{v}_2]$.

\subsection{Coordinate Transformation}

The transformation matrix $\matr{T}$ rotates from global $(x,y)$ to local
$(\bar{x}, \bar{y})$ coordinates. With $c = \cos\theta$, $s = \sin\theta$:
\begin{equation}
    \matr{T} = \begin{bmatrix}
         c & s & 0 & 0 \\
        -s & c & 0 & 0 \\
         0 & 0 & c & s \\
         0 & 0 &-s & c
    \end{bmatrix}
    \label{eq:truss_T}
\end{equation}

\begin{figure}[htbp]
\centering
\begin{tikzpicture}[>=Stealth, scale=1.15]
  %% Element from node 1 (0,0) to node 2 at 35 degrees, length ~4
  \coordinate (N1) at (0,0);
  \coordinate (N2) at ({4*cos(35)},{4*sin(35)});
  %% Bar
  \draw[line width=2.5pt,gray!55] (N1) -- (N2);
  \filldraw (N1) circle (3.5pt) node[below left=1pt] {\small 1};
  \filldraw (N2) circle (3.5pt) node[above right=1pt] {\small 2};
  %% Global axes
  \draw[->,thick] (N1) -- ++(2.2,0) node[right] {$x$};
  \draw[->,thick] (N1) -- ++(0,2.2) node[above] {$y$};
  %% Local axes (rotated by 35 deg)
  \draw[->,thick,blue!75!black] (N1) -- ++({2.0*cos(35)},{2.0*sin(35)}) node[right] {$\bar{x}$};
  \draw[->,thick,blue!75!black] (N1) -- ++({-2.0*sin(35)},{2.0*cos(35)}) node[above left] {$\bar{y}$};
  %% Angle arc and label
  \draw (1.1,0) arc[start angle=0, end angle=35, radius=1.1];
  \node at ({1.35*cos(17)},{1.35*sin(17)}) {$\theta$};
  %% DOF arrows at N2
  \draw[->,dashed,thick,red!65!black] (N2) -- ++(1.0,0) node[right] {\small $u_2$};
  \draw[->,dashed,thick,red!65!black] (N2) -- ++(0,1.0) node[above] {\small $v_2$};
  %% DOF label at N1
  \node[red!65!black,below right=2pt] at (N1) {\small $u_1,\,v_1$};
\end{tikzpicture}
\caption{2D truss element inclined at angle $\theta$: global axes $(x,y)$ and local axes $(\bar{x},\bar{y})$.}
\label{fig:truss_element}
\end{figure}

\subsection{Global Element Stiffness}

The element stiffness in the global frame is:
\begin{equation}
    \matr{k}^e = \matr{T}\transpose \bar{\matr{k}}^e \matr{T}
    = \frac{EA}{L}\begin{bmatrix}
         c^2  &  cs  & -c^2 & -cs  \\
         cs   &  s^2 & -cs  & -s^2 \\
        -c^2  & -cs  &  c^2 &  cs  \\
        -cs   & -s^2 &  cs  &  s^2
    \end{bmatrix}
    \label{eq:truss_global}
\end{equation}

\begin{importantbox}
\textbf{Angle convention:} $\theta$ is measured counter-clockwise from the
positive $x$-axis to the element axis (node 1 $\to$ node 2).
\end{importantbox}
